\chapter{Discussion}
\label{sec:hyper-discussion}
% ===================================================================

\section{Relation to Existing Work}

The book's comparisons with neighbouring programmes --- integrated
information theory, the free energy principle, global workspace,
constructor theory, the ruliad, AIXI, and assembly theory --- are
gathered in Chapter~\ref{ch:related} at the end of the Pledge, where a
reader can find them before committing to six hundred pages rather than
after.
What belongs here is only what is specific to this part.

Turing's oracle hierarchy~\cite{turing1939} is the direct inspiration
for the fibration, and the extension from a linear tower to a fibration
over a category of theories is the departure.
The reinforcement learning framework for sentience is inspired by but
distinct from standard RL~\cite{sutton2018reinforcement}: the reward
signal is \emph{internal} to the account structure rather than
externally provided, and the feeling state is a component of the
resource system that governs computation rather than a scalar to be
maximised.

One correction, since earlier drafts of this chapter made the claim and
it should not stand.
The treatment of consciousness was described here as ordinal-valued and
that description was superseded by Chapters~\ref{ch:apr}
and~\ref{ch:agy}, which locate the vertical coordinate in the Weihrauch
lattice --- emphatically a lattice and not a chain.
Consciousness in this framework is a position an agent occupies, not a
threshold it crosses, and the ordinal tower survives only as the
chain-shaped shadow used to introduce the intuition.

\section{The Broader Picture}

Taken together, the three strands that became this part and its two
neighbours propose a unified framework:

\begin{itemize}[itemsep=4pt]
  \item The physics chapters establish the
        physics: causal structure, dynamics, and quantum amplitudes
        from the bare structure of a GSLT.

  \item The chapters on misled scientists establish the epistemology:
        why situated agents in massive populations perceive a
        hierarchical false ontology rather than the true bisimulation
        quotient.

  \item The present part establishes the phenomenology:
        why complex predictive agents develop graded feeling states
        (sentience) and, with oracle access, inhabit richer worlds
        (consciousness).
\end{itemize}

The unifying theme is that the category of GSLTs is not merely
a setting for studying computation.
It is a setting rich enough to ground physics, epistemology, and
--- if the arguments here have merit --- phenomenology.
The hope is that making these connections mathematically precise
will eventually allow questions about the nature of mind to be
studied with the same rigor currently applied to questions about
the nature of matter.

% ===================================================================
